\section{Stationary gauges and causal modulation}\label{sec:stationary}

Fix one of the horizon-crossing computing cylinders constructed above,
with stationary coordinates $(t,y)$ and exact Kerr metric $b$.
The purpose of this section is to construct a reduced equation whose
stationary zero modes can be controlled by causal parameter feedback.
We first modify the gauge at high spatial frequency, retaining the
stationary spectral information needed at low frequency.  We then
construct the parameter families, the causal target, and the tested
inverse that will supply the compact low-frequency estimate.
All constructions are made separately on each end.  Constants may
depend on the fixed subextremal Kerr parameters and prescribed cylinder.

\subsection{The local spectral input and conventions}

We use the Laplace convention $\partial_t=\zeta$, with
$\operatorname{Re}\zeta\geq0$.  Choose a smooth stationary function
$\phi$ equal to the Kerr tortoise coordinate $r_*$ near infinity and
zero on a fixed interior collar.  Outgoing fields have the factor
$e^{-\zeta\phi}$.  At infinity our convention is
\begin{equation}\label{st:b-weights}
 H_{\mathrm b}^{s,\ell}
 =r^{-\ell-3/2}H_{\mathrm b}^{s}(dr\,d\omega/r),
\end{equation}
where the derivatives in $H_{\mathrm b}^{s}$ are $r\partial_r$ and
unscaled angular derivatives.  Interior spaces are extendible at the
strict inner outflow boundary; their dual spaces are supported across
that boundary.  We denote these conventions by a bar and a dot,
respectively.  The notation $O(r^{-j+})$ permits an arbitrarily small
loss in the exponent, chosen within the relevant indicial interval.

\begin{proposition}[Local spectral input]\label{st:hhv}
On the fixed subextremal Kerr cylinder there are arbitrarily small,
stationary, smooth, compactly based bundle-map modifications of
harmonic gauge and of its symmetric-gradient partner with these
properties.
\begin{enumerate}
\item The stationary gauge-potential operator, obtained by applying
the modified gauge to a pure Lie derivative, is invertible at decay
strictly between the powers $1$ and $r^{-1}$.
\item The modified linearized reduced vacuum operator has outgoing
inversion at nonzero frequency in the closed physical half-plane.
Its bounded-frequency estimates, with a weaker compact error, are
uniform down to zero in the weight interval $-3/2<\ell<-1/2$.
Its zero-frequency realization has index zero and a seven-dimensional
kernel: four gauge-corrected Kerr variations and three translations.
\item In outgoing time, a polynomial mode with coefficients of
order $O(r^{-1+})$ has translation leading coefficient if its
degree is one.  There are no such modes of degree at least two.
\end{enumerate}
For tensors one may take $s>5/2$, with strict margin, also sufficient
for gauge potentials.  The uniform near-zero equation norm is
$H_{\mathrm b}^{s-1,\ell+2}$.  At nonzero frequency the outgoing
conditions are $\ell<-1/2$ and $s+\ell>-1/2$.
\end{proposition}

\begin{proof}
The gauge-potential assertion is \cite[Theorem~5.1]{HHV2025};
the discussion following that theorem permits arbitrarily small local
damping parameters.  The bounded-frequency estimates and stationary
index are \cite[Theorem~5.2]{HHV2025}, and nonzero-frequency inversion
and the seven-dimensional kernel are \cite[Theorem~6.6]{HHV2025}.
The polynomial-mode assertions are
\cite[Proposition~6.8 and Lemma~6.9]{HHV2025}.
We verify the geometric conventions for applying these results here.
The stationary time is horizon-regular, with outgoing time
$t-\phi$ at infinity.  In Boyer--Lindquist coordinates there,
\[
 r_*=r+2M\log r+O(1),\quad b^{ti}=O(r^{-2}),\quad
 \partial_y(\phi-r)=O(r^{-1}),\quad
 b^{-1}\bigl(d(t-\phi),d(t-\phi)\bigr)=O(r^{-1}),
\]
with symbol derivatives.  The logarithmic coefficient is independent
of spin.  Bounded smooth order-zero offsets in $\phi$ preserve the
outgoing graphs and polynomial-mode statements.  Stationary regular
chart changes on compact slabs act by smooth component maps and
bounded time shifts, with only polynomial losses on fixed strips.

The inner edge may be extended to the chosen $r_b$ between the
horizons.  Between any two strict depths, $dr$ is temporal, so
ordinary hyperbolic propagation applies.  On a strip with bounded
$\operatorname{Re}\zeta$, its high-frequency cost is polynomial:
$\operatorname{Im}\zeta$ is a real tangential Fourier parameter in
the radial-time Cauchy estimate, and the real Laplace part is a
bounded lower-order weight.  We choose the local damping sufficiently
small for the finitely many strict thresholds below.  This input
concerns the local modified operator; it invokes neither nonlinear
stability nor inversion after spatially nonlocal gauge changes.
\end{proof}

Normalize the reduced equation by
\begin{equation}\label{st:reduced-normalization}
 \operatorname{Ric}(g)+(\mathcal D_g+C)\Upsilon=0,\qquad
 (\mathcal D_g\omega)_{\mu\nu}=\nabla^g_{(\mu}\omega_{\nu)}.
\end{equation}
For a target metric $B$, the differential wave-map gauge is
\begin{equation}\label{st:wave-gauge}
 (\Upsilon_{\mathrm{wave}}(g;B))_\mu
 =-g_{\mu\kappa}g^{\alpha\beta}
   \bigl(\Gamma(g)^\kappa_{\alpha\beta}
              -\Gamma(B)^\kappa_{\alpha\beta}\bigr).
\end{equation}
Write $E_{\mathrm g}$ for the addition to this gauge and $C$ for
the addition to its partner, including the local terms just obtained.
At $g=B=b$ write
\[
 \mathcal B=D\Upsilon_{\mathrm{wave}}|_{b;b}+E_{\mathrm g},
 \qquad
 L=\operatorname{Ric}'_b+(\mathcal D_b+C)\mathcal B.
\]
The notation $L(\zeta)$ means substitution of $\partial_t=\zeta$.
We write $L_{\mathrm h}$ for the local operator before the
high-spatial-frequency changes made next.

\subsection{Gauge transport on the spatial elliptic cone}

The geometric analysis supplies a compactly based spatial conic
neighborhood $\mathcal T$ of both projected trapped sets, disjoint
from the horizons and ends of the cylinder, on which the
zero-frequency spatial principal symbol is strictly elliptic.
We use this cone instead of a timelike neighborhood in the spatial
base.  The two temporal characteristic roots remain separated because
the $t$-slices are spacelike.

For this construction let $\bar g$ be an exact nearby stationary Kerr
metric in a chart of the same type, and let $B_0$ be a close stationary
target.  We calibrate the transport of
\begin{equation}\label{st:frozen-general}
 \operatorname{Ric}'_{\bar g}
 +(\mathcal D_{\bar g}+C)
       \bigl(D\Upsilon_{\mathrm{wave}}|_{\bar g;B_0}
                         +E_{\mathrm g}\bigr).
\end{equation}
At equal metrics without additions, this is ordinary tensor parallel
transport after removal of scalar wave-amplitude transport.
Changing $B_0$ inserts only an algebraic term $K_{\mathrm{alg}}$
in the linearized gauge. We use these wave-gauge identities in the
following construction.

\begin{lemma}[Conic gauge calibration]\label{st:conic-gauge}
The additional gauge and partner terms can be chosen as real spatial
pseudodifferential operators $A_0+A_1\partial_t$, with respective
spatial orders $0,-1$, such that:
\begin{enumerate}
\item their kernels are compactly supported in both variables away
from horizons and boundaries, and their high-frequency microsupport
lies in $\mathcal T$;
\item at sufficiently high spatial frequency their transport near
each trapped set is exactly compatible with a smooth positive
homogeneous bundle metric;
\item at the base metric the unsaturated transport and cutoff
transition have imaginary subprincipal defect smaller than any
prescribed fixed fraction of the minimum normal expansion rate;
\item the additional temporal-acceleration coefficients have spatial
order at most $-1$ and are small on any prescribed finite Sobolev
list when the high-frequency cutoff $\Lambda$ is sufficiently large.
\end{enumerate}
The prescription is smooth in the stationary metric and target.
The symbol bounds are uniform for large $\Lambda$ and for a
sufficiently small fixed parameter neighborhood.
\end{lemma}

\begin{proof}
For a null covector $k$ define
\[
 Q_kH=H(k^\sharp,\cdot)
                  -\tfrac12(\operatorname{tr}_{\bar g}H)k,
 \qquad d_ke=k\odot e,
\]
where $\odot$ is averaged symmetric product.  The map $d_k$ is
injective and $Q_k$ is surjective, since trace reversal is invertible.
Nullity gives $\operatorname{im}d_k\subset\ker Q_k$.
Parallel transport preserves this flag, with quotient the traceless
symmetric tensors of the Euclidean screen.  Up to a common scalar
normalization the available changes are $d_kE_{\mathrm g}(k)$
and $C(k)Q_k$.

Normalize $k$ by its conserved nonzero stationary energy:
$n=k^{\sharp_{\bar g}}/E$ is parallel along its null geodesic.
Choose a stationary future $\bar g$-unit timelike field $T$ and set
\[
 \alpha=-\bar g(n,T)>0,\qquad
 l=\frac{T}{\alpha}-\frac{n}{2\alpha^2}.
\]
Thus $\bar g(l,l)=0$ and $\bar g(n,l)=-1$.  In the screen
$\{n,l\}^{\perp}$ choose local orthonormal vectors $e_A$.
Writing $\nabla$ for transport along the geodesic, preservation
of these inner products gives
\[
 \nabla n=0,\qquad \nabla l=w,\qquad
 \nabla e_A=\sum_B\omega_{AB}e_B+\bar g(w,e_A)n,
 \qquad \omega_{AB}=-\omega_{BA},
\]
where $w$ is a screen vector.  In particular there is no diagonal
boost term: conservation of $E$ has made $n$ parallel, not merely
parallel up to scale.  Give $n$ positive length $\epsilon$, $l$
length $\epsilon^{-1}$, and the screen its ordinary metric.
The null-rotation terms have norm $O(\epsilon)$ in this metric,
whereas the screen rotation is skew.  The construction patches
without a global screen frame and gives a positive homogeneous
tensor metric.  Its non-skew transport is arbitrarily small,
also on the fixed tensor power.

Choose metric transport preserving the orthogonal flag splitting,
with the same induced quotient transport.  Its difference $A$
from the original transport maps $\ker Q_k$ into
$\operatorname{im}d_k$ and is otherwise unrestricted.  It is
therefore realized by the available changes: for a smooth right
inverse $R_k$ of $Q_k$, put $P_k=1-R_kQ_k$ and take
\[
 E_{\mathrm g}(k)=d_k^{-1}AP_k,\qquad C(k)=AR_k.
\]
Then $d_kE_{\mathrm g}(k)+C(k)Q_k=A$.  Compensate also
$K_{\mathrm{alg}}$ and the local additions in prescribing the
total transport.  Smooth projections and right inverses give
smooth symbols.  The normalized-ray construction can be chosen
real and invariant under simultaneous covector reversal.
Transport before time normalization is odd of degree one,
and the gauge-factor symbols are even of degree zero.

At $b$ choose the local additions small after the flag scaling.
Convex interpolation to compatible transport preserves the strict
subprincipal margin through the high-frequency transition.
Write the roots as $k=(-P_\pm(y,\eta),\eta)$ and let
$E_\pm^{\mathrm{add}}$ be the extra value prescribed for either
gauge factor on the two roots.  Set
\begin{equation}\label{st:root-interpolation}
 A_1=\frac{E_+^{\mathrm{add}}-E_-^{\mathrm{add}}}
                  {-i(P_+-P_-)},\qquad
 A_0=E_+^{\mathrm{add}}+iP_+A_1.
\end{equation}
Since $|P_+-P_-|\gtrsim|\eta|$, these have orders $-1,0$.
Reversal interchanges the sheets; thus $A_0$ is real-even,
$A_1$ imaginary-odd, and $A_0+A_1\partial_t$ is real.
Extend and patch the values on each sheet before the single
spatial interpolation.  Overlap of their projections creates no
conflict, including over the ergoregion.

Use an even conic cutoff in $\mathcal T$, equal to one near both
projected trapped sets, and a cutoff $|\eta|\gtrsim\Lambda$.
Quantize with compact kernel cutoffs in both variables.  Products
using both order-zero gauge factors have no remaining first
derivative and do not contribute to transport order.  Thus the
common scalar cutoff really interpolates the first-order transport;
symbol derivatives in compositions are lower order.  Kernel
localization preserves its homogeneous symbol.  Outside a slightly
enlarged cone, separated-support composition or integration by
parts away from the spatial diagonal gives smoothing leakage
$O(\Lambda^{-N})$ in every prescribed fixed seminorm, for every
fixed $N$.  The full symbols are uniformly bounded in their
orders.  Each extra coefficient of a second time derivative has
a spatial gain of at least one, hence is small on any fixed
finite Sobolev list.  The mode-preserving finite-rank corrections
below are smoothing and do not alter the calibration.  Repeating
the prescription at each frozen metric makes compatibility exact
near its trap.
\end{proof}

\subsection{Stationary estimates after the conic modification}

We separate bounded frequency, high frequency, and the outgoing
spatial transpose.  The last is the bilinear \emph{spatial}
transpose, whose outgoing exponential has the same sign as the
primal exponential. Fix a small weight slack $0<\iota\ll1$ in the
weighted scattering estimates below.

\begin{proposition}[Stationary estimates]\label{st:stationary-estimates}
For sufficiently large $\Lambda$, the conically modified base
operator has the following estimates.
\begin{enumerate}
\item Bounded-frequency estimates with weaker compact errors persist
uniformly as $\Lambda\to\infty$, at the orders and zero-frequency
weights of Proposition~\ref{st:hhv}.  The stationary index is zero,
and kernel and cokernel are upper semicontinuous.  On closed bounded
bands excluding zero, the physical estimates persist without compact
error, up to $\operatorname{Re}\zeta\downarrow0$.
\item For every fixed strip $0\leq\operatorname{Re}\zeta\leq D_*$
and fixed admissible order, there are polynomial high-frequency
bounds without kernel, with threshold independent of sufficiently
large $\Lambda$.  These include the compact $H^s$ estimate from
$H^{s-1}$, for $s=5$ or a fixed higher admissible order, the supported
dual solution estimate of order $1-s$, and
\[
 \langle r\rangle^{-1/2-\iota}(v,\partial_yv)\in L^2
\]
for the physical transpose with a compact smooth source, with
polynomial frequency cost.
\item Zero-frequency estimates, including those for gauge
potentials, persist uniformly for nearby exact Kerr principal
metrics in smooth stationary, asymptotically affine charts of
the prescribed type.  Fixed smooth finite-dimensional supplements
remove their kernel and cokernel locally in parameters.
\end{enumerate}
The same-sign outgoing transpose has the a priori estimate and
index stated in Lemma~\ref{st:transpose}.
\end{proposition}

\begin{proof}
\emph{Bounded frequency.}
The additional operator has spatial order at most one and compact
elliptic microsupport.  Let $\beta$ be an order-zero elliptically
supported cutoff equal to one on an enlarged neighborhood of that
microsupport.  On fixed enlarged compact sets $K'$,
\begin{equation}\label{st:elliptic-smallness}
 \|(L-L_{\mathrm h})(\zeta)u\|_{H^{s-1}}
 \leq C\Lambda^{-1}
       \|\operatorname{Op}(\beta)u\|_{H^{s+1}}
       +O(\Lambda^{-N})\|u\|_{H^s(K')}.
\end{equation}
Nested cutoffs and the high-frequency cutoff give this estimate,
including arbitrary-order separated-symbol remainders.  An elliptic
parametrix for the operator without additions controls the stronger
norm.  The addition is already bounded $H^s\to H^{s-1}$, so the
bounded-frequency interior graph domains agree.  The same argument
applies to supported dual orders, since reversed covectors remain
elliptic, and to gauge potentials.  It proves persistence of
compact-error estimates, index, and semicontinuity.  Smoothing
subtractions of size $O(\Lambda^{-N})$ do not change this reasoning.

We recall why physical homogeneous limiting-absorption states on
a closed nonzero band belong to the outgoing class in
Proposition~\ref{st:hhv}.  The passage from physical limiting
absorption to outgoing conormality uses the unchanged Kerr end and
the interior propagation verified above.  The physical solution and
source weights are $\langle y\rangle^{-1/2-\iota}$ and
$\langle y\rangle^{1/2+\iota}$.  The incoming radial estimate gives
the outgoing microlocal condition for a homogeneous real-frequency
limit.  Exterior rotation commutators are short-range because the
principal $1/r$ correction is rotation invariant.  Divided rotation
differences, exterior propagation, and spatial ellipticity for
unscaled derivatives and fixed-annulus errors give the same
scattering estimate after any fixed number of rotations, with
positive slack.  Thus, for a limit $h_{\mathrm{lim}}$ of the local
operator, $rh_{\mathrm{lim}}$ and its radial derivative lie in
$L^2(r^{-1-2\iota}dr\,d\omega)$ with those rotations.

Divide the radial equation by its exact radial coefficient.
Its leading factorization in tortoise coordinate is
$(\partial_{r_*}-\zeta)(\partial_{r_*}+\zeta)(rh_{\mathrm{lim}})$;
the remaining symbol terms are $O(r^{-2})$ with at most two
angular derivatives or one radial and one angular derivative.
Integrating the first factor from infinity, with the incoming
homogeneous term excluded, gives $O(r^{-1+\iota})$ for the
second factor by weighted Cauchy--Schwarz.  On the imaginary
axis put $w=e^{\zeta r_*}rh_{\mathrm{lim}}$.  It has size
$O(r^\iota)$ and the first symbol-derivative estimate, and
\begin{equation}\label{st:outgoing-symbol}
 (\partial_{r_*}-2\zeta)\partial_{r_*}w
 =O(r^{-2})\bigl(\partial_\omega^{\leq2}w,
                  \partial_\omega^{\leq1}\partial_{r_*}w\bigr).
\end{equation}
Here is the derivative induction explicitly.  Write the right
side as $F=r^{-2}(A_2w+A_1w')$, with symbol-bounded angular
operators $A_2,A_1$ of orders at most two and one.  Initially
$w=O(r^\iota)$ and $w'=O(r^{-1+\iota})$.  The equation first
gives the raw bound $w''=2\zeta w'+F=O(r^{-1+\iota})$.
It follows that $F'=O(r^{-3+\iota})$, including the term
$r^{-2}A_1w''$.  Differentiating the equation and excluding its
nondecaying homogeneous solution gives
\[
 w''(r_*)=-\int_{r_*}^{\infty}
       e^{2\zeta(r_*-s)}F'(s)\,ds=O(r^{-2+\iota}).
\]
In general, symbol bounds through $w^{(j-1)}$ first give the raw
bound $w^{(j)}=O(r^{-j+1+\iota})$.  Then
$F^{(j-1)}=O(r^{-j-1+\iota})$, including its highest derivative
term, and integration of
$(\partial_{r_*}-2\zeta)w^{(j)}=F^{(j-1)}$
gives $w^{(j)}=O(r^{-j+\iota})$.  The raw decay excludes the
homogeneous term also when $\operatorname{Re}\zeta=0$.
Angular derivatives are included throughout, and initial slack
is chosen smaller than the intended final slack.  This gives outgoing b-regularity,
and radial propagation upgrades interior regularity above threshold.
Proposition~\ref{st:hhv} excludes the resulting homogeneous limit.
In the open half-plane one may instead use its outgoing primal
solvability and the physical index argument below.
The microlocal smallness \eqref{st:elliptic-smallness} now retains
the closed-band estimates for $L$.  Dual estimates follow from
the open half-plane.  The contour argument will use bounds up
to the imaginary axis, not an ordinary unweighted inverse there.

\emph{High frequency.}
We apply the local outgoing normally hyperbolic estimate of
\cite[Theorem~1]{Dyatlov2016}, with the scalar-principal bundle
formulation and positive pseudodifferential metrics of
\cite[Condition~(2.5), Proposition~3.12]{Hintz2017}.  Its hypotheses are a symplectically
transverse codimension-two trapped set, smooth stable and unstable
surfaces, and imaginary subprincipal transport defect strictly
below half the minimum normal rate in a positive metric, including
its transport and the half-density convention.  The geometric
defining functions $z_1\pm z_2$ and Lemma~\ref{st:conic-gauge}
verify these hypotheses for each nonzero-energy sign with margin.

To use the separated symbol $q_E=\Sigma_Kb^{-1}(k,k)$ of
\eqref{geo:separated-symbol}, an operator initially written with principal
symbol $b^{-1}(k,k)$ can be normalized
microlocally by symmetric sandwiching with the smooth elliptic
multiplier $\sqrt{\Sigma_K}$.  In half-density Weyl calculus the
two first-order bracket terms in this symmetric sandwich cancel,
so its subprincipal symbol is $\Sigma_K$ times the original one.
On the characteristic trap $H_{\Sigma_K p}=\Sigma_KH_p$;
the derivative of a chosen bundle metric scales by the same
factor.  Consequently the pointwise relative transport margin
is preserved.  For a theorem stated with the supremum of the
defect and the minimum normal rate, fix this normalization first
and choose the defect sufficiently small using the positive
minimum and finite maximum of the scaling factor on the compact
trapped neighborhood.  The gauge calibration permits
arbitrarily small defect before saturation and zero defect at
saturation, so this choice imposes no additional hypothesis.

For $k=-E\,dt+\eta$,
\[
 \partial_E\bigl(b^{-1}(k,k)\bigr)=-2k^\sharp(dt).
\]
Its sign is separated on each characteristic sheet because $dt$
is timelike, also after rescaling at spatial fiber infinity.
Orient projected propagation by increasing regular $t$, dividing
the Hamilton field by its $t$-speed.  At trapping
$k^\sharp(dt)/E>0$; the frequency and time normalizations differ
by a bounded positive factor after fixing orientations.  The
subprincipal margin is selected after these normalizations.
More precisely, if $f$ is the positive time-change factor,
$f_{\max}$ its maximum on the compact microlocal neighborhood,
and $\nu_{\min}^{\mathrm{new}}>0$ the minimum normal rate for the
normalized flow, choose the arbitrarily small flag defect
$\delta$ so that
$f_{\max}\delta<\nu_{\min}^{\mathrm{new}}/2$ with fixed surplus.
Choose the local additions smaller than that surplus.
Compatible saturated transport has zero defect, and the convex
cutoff transition retains the margin.  Merely scaling rates
and defects pointwise would not suffice for this global
supremum--infimum comparison.  Scalar half-density transport is scalar
wave-amplitude transport with stationary density; its remaining
bundle transport is the calibrated one.  The bounded positive real
Laplace part therefore supplies damping on either trapped sheet.
The horizon radial operator is unchanged and has positive redshift.

The ray alternatives give strict outflow and no returns from
infinity or homoclinic returns.  Incoming propagation from infinity,
ordinary real-principal propagation, radial estimates, and the trap
estimate control the characteristic set modulo lower remainders.
For the transpose use the opposite propagation direction.
A fixed Killing-energy sign can encounter ergoregion rays with
a different time orientation.  They cannot be trapped, since
the double-root identities with future $P>0$ force future $E>0$.
They use only escape or radial propagation.  At spatial fiber
infinity with normalized temporal energy zero, the new symbols
miss all nonelliptic directions and the horizon radial geometry;
the zero-energy alternatives give ordinary nonradial propagation.
The positive metric extends to the required microlocal neighborhood.
This uses a local outgoing theorem, not global Kerr--de Sitter
inversion or a positive Killing-current norm in the ergoregion.

We check uniformity in $\Lambda$.  On characteristic added
microsupport, $c|\zeta|\leq|\eta|\leq C|\zeta|$.  Put
$h_f=|\zeta|^{-1}$ and $\xi=h_f\eta$.  On
$|\eta|\geq c'|\zeta|$, the additional terms in $h_f^2L$ are
uniformly $O(h_f)$ of maximal spatial order one, by the ordinary
$S^j_{1,0}$ bounds on coefficients of orders $1,0,-1$ multiplying
zero, one, and two powers of $\zeta$.  Outside an enlarged elliptic
cone they are $O(h_f^\infty)$ smoothing.  Where the cutoff changes
at the trap, $\Lambda/|\zeta|$ is bounded above and below, so its
rescaled derivatives are uniformly bounded.  Its first transport
is the controlled scalar interpolation; symbol derivatives in
gauge composition are lower order.  Substitution of the bounded
real Laplace part in the extra time-derivative factors also changes
only lower orders.  At large relative spatial frequency the
added microsupport is itself elliptic, and the additions vanish
at the radial set.

For $|\eta|\leq c'|\zeta|$, with $c'$ sufficiently small, the wave
symbol has size $\gtrsim|\zeta|^2$.  Write the perturbation as
\[
 \zeta^2K_{-1}+\zeta K_0+K_1,
\]
with subscripts denoting maximal spatial orders.  Relative to
elliptic inversion its first term costs $O(\Lambda^{-1})$ and
the others $O(|\zeta|^{-1})$.  Separated localizations from this
sector to the characteristic set have rapid frequency decay,
including for fixed-rank smoothing terms.  No uniform smooth
semiclassical symbol at $\xi=0$ is asserted.  These errors are
absorbed after propagation with a threshold independent of large
$\Lambda$.  Large positive Laplace lines also admit forward wave
energy after the small acceleration correction is solved.
This proves the polynomial strip bounds and the exterior
scattering bounds.

\emph{Nearby stationary parameters.}
At zero frequency the horizon inequalities remain strict after
pullback to nearby exact Kerr principal models.  Lower-order terms
vary smoothly, high additions remain elliptic and away from the
horizon, and the normal operator at infinity is a positive
constant-coefficient Laplacian after a constant linear spatial
change, with unchanged indicial gaps.  Add a fixed smooth
finite-dimensional supplement removing the center's kernel and
cokernel.  Primal and dual compact-error estimates, compactness,
and equation convergence keep the supplemented problem invertible:
otherwise a normalized contradicting sequence yields a nonzero
limiting nullstate.  This avoids demanding smallness of a changing
principal operator on a fixed one-derivative graph.  Gauge
potentials obey the same argument.  When high spatial regularity is
needed only on a compact set, the b-elliptic end permits lower
solution and equation derivative orders after a switch across
fixed full elliptic annuli, with the higher source norm maintained
slightly farther out.
\end{proof}

The transpose argument uses three realizations of the stationary
operator.  The physical equations act on the original fields through
$L(\zeta)$ and its supported spatial transpose $L(\zeta)^t$.
Removing the outgoing exponential gives, respectively,
\[
 P_{\rm out}(\zeta)=e^{\zeta\phi}L(\zeta)e^{-\zeta\phi},
 \qquad
 P^\sharp(\zeta)=e^{\zeta\phi}L(\zeta)^t e^{-\zeta\phi}.
\]
The second is a same-sign outgoing conjugation of the physical
transpose.  In particular,
\[
 P_{\rm out}(\zeta)^t=e^{-\zeta\phi}L(\zeta)^t e^{\zeta\phi}.
\]
The conjugating factors are reversed when the outgoing primal is
transposed.  We first estimate $P^\sharp$ near zero and then compute
its index at nonzero frequency, using different exterior equation
spaces for these two purposes.

\begin{lemma}[The outgoing spatial transpose]\label{st:transpose}
Fix $\ell_{\mathrm{wt}}=-1/2-\iota$, with $0<\iota\ll1$, and
a sufficiently large fixed exterior derivative order $j$.
Let $\mathcal A$ have supported interior order $-4$ and exterior
space $H_{\mathrm b}^{j,\ell_{\mathrm{wt}}}$; let $\mathcal S$
have supported interior order $-5$ and exterior space
$H_{\mathrm b}^{j-1,\ell_{\mathrm{wt}}+2}$.
For
\[
 P^\sharp(\zeta)=e^{\zeta\phi}L(\zeta)^t e^{-\zeta\phi}
\]
there is an estimate from $\mathcal S$ to $\mathcal A$ with a
weaker compact error, uniform near zero.
At each fixed small nonzero $\zeta$ with
$\operatorname{Re}\zeta>0$, the graph using instead the exterior
equation space $H_{\mathrm b}^{j,\ell_{\mathrm{wt}}+1}$
has index zero, also after compact smooth finite-rank augmentation.
\end{lemma}

\begin{proof}
Combine the supported interior stationary dual estimate with the
primal outgoing estimate cut to infinity.  In Cartesian
trivializations the difference of the conjugated transpose and
comparison primal operator is
\[
 O(r^{-2})\partial_y+O(r^{-3})
                     +(\zeta+\zeta^2)O(r^{-2}).
\]
The symbol bounds, notably $b^{ti}=O(r^{-2})$ and connection
coefficients of that order, give this expression.  It is small
between the indicated spaces sufficiently far out and near zero.
Switch orders across fixed full elliptic annuli beyond all
compact kernels.  This proves the a priori estimate.

For the index, first the
unconjugated physical primal and supported transpose have index
zero for each $\operatorname{Re}\zeta>0$.  Use interior
solution/equation orders $s,s-1$ with $s=5$ and $1-s,-s$,
respectively, and ordinary unweighted fixed-frequency graphs
outside.  The Fourier normal operator at infinity is a multiple
of the flat spatial Laplacian minus $\zeta^2$, hence invertible.
At high interior spatial frequency the normalized temporal
energy is zero; ordinary propagation, non-adverse redshift, their
supported dual estimates, and strict outflow give compact-error
estimates in both directions.  The additions are lower order over
ellipticity or smoothing with compact kernels.  Thus the graph
domains are locally constant in the open half-plane and both
problems are Fredholm.

On a sufficiently large forward line, wave energy and past
integration give primal inversion and uniqueness after solving
the small order-$(-1)$ normal-acceleration correction.  Begin
with natural unweighted orders, then use exterior ellipticity.
Duality gives transpose solvability; testing against primal
solutions at slightly raised compact orders gives its uniqueness.
Continuation establishes the physical indices, which compact
smooth finite-rank terms do not change.

The comparison outgoing primal problem uses $L_{\mathrm h}$ with
extendible interior orders and exactly the exterior graph proposed
for the outgoing transpose.  Proposition~\ref{st:hhv} gives
solvability on compact smooth data, trivial kernel, and an estimate
with compact error.  Altering interior orders uses the one-order
propagation estimate, homogeneous regularity above threshold,
and annular ellipticity.  This comparison also has index zero.
Exchange its end with the physical supported transpose end in
their direct sum, by a smooth signed exchange on an elliptic
annulus outside all kernels.  Domain and range spaces identify,
and collar errors are Rellich compact.

On the conjugated end the fixed-frequency difference is
$O(r^{-2})\partial_y+O(r^{-2})$.  Its graph additionally controls
$H_{\mathrm b}^{j+1,\ell_{\mathrm{wt}}-1}$.  Indeed, if
$V\in H_{\mathrm b}^{j,\ell_{\mathrm{wt}}}$ and its equation
$PV\in H_{\mathrm b}^{j,\ell_{\mathrm{wt}}+1}$, write
\[
 r^2P=B_2+rB_1+B_0,
\]
where $B_2$ is b-elliptic of order two and the remaining terms
have order at most one, with uniformly bounded symbol
coefficients at this fixed nonzero frequency.  Both $r^2PV$
and $(rB_1+B_0)V$ lie in
$H_{\mathrm b}^{j-1,\ell_{\mathrm{wt}}-1}$; elliptic regularity
gives the asserted extra derivative.  Since
$r^{-2}\partial_y$ is $r^{-3}$ times a b-vector field, the
first-order difference maps this improved graph regularity to
$H_{\mathrm b}^{j,\ell_{\mathrm{wt}}+2}$.  The zeroth-order
difference maps the original $V$ to the same space.
Thus the difference gains one tail weight relative to the
equation space $H_{\mathrm b}^{j,\ell_{\mathrm{wt}}+1}$, with
norm $O(Y^{-1})$ beyond radius $Y$.  On each finite annulus
ordinary graph ellipticity gives two local derivatives over
the equation order, so Rellich compactness handles the
remaining compact part.
The other exchanged summand is a physical primal problem of
index zero.  Hence the desired outgoing transpose graph has
index zero.  All high additions remain on their original interior
pieces in this argument.

This gives an index, not automatically an inverse.  If a compactly
forced outgoing conjugated transpose problem has trivial kernel,
index zero supplies its inverse.  Its solution $V$ gives the
physical test $v=e^{-\zeta\phi}V$, exponentially decaying at
infinity.  Supported--extendible integration by parts against
these tests kills the physical primal kernel, and physical index
zero yields inversion.
The stronger near-zero $\mathcal S$ estimate is used on graph
states with that extra source weight.  This never identifies
$P^\sharp(\zeta)$ with the transpose of the conjugated primal.
\end{proof}

\subsection{Stationary profiles and exact parameter families}

We next retain the physical zero modes while installing the conic
gauge, and integrate the relevant directions into exact Kerr
families. The indicial calculation below identifies the required
parameter directions at the fixed center.

\begin{lemma}[Profiles and mode-preserving corrections]\label{st:profiles}
There are smooth stationary tensors
$H_i$ for $1\leq i\leq4$ and $J_i,H_i^1$ for $1\leq i\leq3$,
with stationary vector fields $V_i,W_i$, such that
\begin{gather}
 J_i=\mathcal L_{V_i}b,\qquad
 H_i^1=\mathcal L_{W_i}b+2dt\odot V_i^\flat,
 \label{st:profile-definition}\\
 V_i=e_i+O(r^{-1}),\qquad
 W_i=y_i\partial_t+O(r^{0+}),
 \label{st:vector-profiles}\\
 \mathcal BH_i=\mathcal BJ_i
   =\mathcal B(tJ_i+H_i^1)=0,
 \qquad
 L(0)H_i=L(0)J_i=0,\qquad
 L(0)H_i^1=-[L,t]J_i.
 \label{st:profile-identities}
\end{gather}
Here $e_i$ are coordinate translations near infinity and $H_i$
are gauge-corrected Kerr variations.  In symbol orders,
\begin{equation}\label{st:profile-tails}
 H_i,H_i^1=r^{-1}(\text{smooth angular term})+O(r^{-2+}),
 \qquad
 J_i=r^{-2}(\text{smooth angular term})+O(r^{-3+}).
\end{equation}
The additional high-frequency gauge can be changed by
$O(\Lambda^{-N})$ finite-rank smoothing terms, for every fixed
$N$, so that these identities remain exact.  Its decaying
stationary kernel then consists precisely of the seven
independent $H_i,J_i$.
\end{lemma}

\begin{proof}
For the local gauge, the starting profiles are the Kerr, translation,
and boost profiles of \cite[Lemmas~6.1--6.3, Proposition~6.5,
and~(6.6)]{HHV2025}.  We record the end estimates needed here, including
the sharper translation-potential bound, and then preserve their
identities under the conic modification.

Begin with the local gauge of Proposition~\ref{st:hhv}.
Correct cut-off Kerr, translation, and boost generators
successively with its stationary gauge-potential inverse.
The errors requiring nondecaying correction, for Kerr and the
stationary boost part, are $O(r^{-2})$.  Shifting to a slight
growth weight $O(r^{0+})$ is surjective by indicial Fredholm
regularity: its cokernel is already zero at decay.

The translation correction has an $O(r^{-3})$ leading error
which is a constant linear combination of second Cartesian
derivatives of degree-$(-1)$ symbol-leading coefficients.
Each has zero spherical mean, as follows by integration over
shells and surface scaling.  Hence the decaying vector correction
has an $r^{-1}$ angular leading term without a leading logarithm.
The identity
$\mathcal L_{te_i+y_i\partial_t}\eta_{\mathrm{Mink}}=0$
and the gauge correction give
$tJ_i+H_i^1=\mathcal L_{tV_i+W_i}b$ and
\eqref{st:profile-identities} for the local gauge.
The flat-Laplacian indicial expansion gives
\eqref{st:profile-tails}, using $b^{ti}=O(r^{-2})$ in the
equation for $H_i^1$.
The sublinear correction vectors for Kerr and for
$W_i-y_i\partial_t$ have first derivatives $O(r^{-1})$
with symbol bounds.  Indeed their flat symmetric gradients
have that order by the tensor profile expansion.  Differentiating
and permuting indices gives second derivatives of order $r^{-2}$;
integrating along rays from infinity, where first derivatives
vanish, gives the stated bound.  These arguments use smooth Kerr
symbols, indicial regularity at infinity, and compact propagation
above threshold, without a restriction on spin.

The seven stationary profiles are independent already on a
sufficiently far exterior region, away from the trapping base.
A stationary sublinear-symbol gauge preserves the $1/r$
coefficient of the stationary Killing-time norm and the leading
curl of $b_{ti}\,dy_i/b_{tt}$.  A stationary time displacement
adds an exact form to the latter, and its spatial Lie derivative
is lower order.  The Kerr expansions therefore detect the mass
and angular momentum vector, with fixed normalized Cartesian
axes.  A nonzero translation combination cannot have zero Lie
derivative there, since it commutes with $\partial_t$ and its
action on the leading radial coefficient of
$b(\partial_t,\partial_t)$ is nonzero.  Thus compact smooth
functionals away from a physical trapping neighborhood separate
the modes; no elliptic unique continuation across that
neighborhood is required.

At the base point, subtract finite-rank spatial smoothing terms
from the high addition to $E_{\mathrm g}$ so that its $A_0$
part annihilates the span of $H,J,H^1$ and its $A_1$ part
annihilates $J$.  To be explicit, choose a compact cutoff
$\chi_{\mathrm{in}}$ equal to one on all input supports.
Let $q_\alpha$ be a basis of the actual finite-dimensional span
of $\chi_{\mathrm{in}}H,\chi_{\mathrm{in}}J,
\chi_{\mathrm{in}}H^1$, and choose compact smooth tests
$\lambda_\alpha$ with $\lambda_\alpha(q_\beta)=\delta_{\alpha\beta}$.
For the high-cutoff operator $A_{0,\Lambda}$ subtract
\[
 S_0f=\sum_\alpha (A_{0,\Lambda}q_\alpha)
                         \lambda_\alpha(\chi_{\mathrm{in}}f).
\]
The identity $A_{0,\Lambda}=A_{0,\Lambda}\chi_{\mathrm{in}}$
shows that $A_{0,\Lambda}-S_0$ annihilates every original
profile, including when their restrictions are dependent.
Use the same construction on the span of
$\chi_{\mathrm{in}}J$ for $A_{1,\Lambda}$.
Each $q_\alpha$ is fixed, smooth, and compactly supported.
Consequently $A_{0,\Lambda}q_\alpha=O(\Lambda^{-N})$
in each fixed smooth seminorm for every fixed $N$, and the
finite-rank kernel formula gives the same statement for $S_0$
and $S_1$ in every fixed smoothing norm.  The constants may
depend on that norm and on $N$: only the finite base list is
made small by the single choice of $\Lambda$.  Hence the
profile identities remain exact.  The seven modes persist,
and Proposition~\ref{st:stationary-estimates} excludes any
enlargement of the decaying kernel.  Gauge-potential
invertibility persists by the same elliptic smallness.
Index and weak compactness give the corresponding cokernel
continuity towards the local-gauge limit.  Extend these
smoothing subtractions smoothly to nearby parameters.
\end{proof}

The decaying stationary modes are not the whole kernel when a field
is allowed slight growth at infinity.  A constant symmetric covariant
two-tensor in four dimensions has ten independent coefficients.
The next proposition realizes these constant asymptotes by additional
homogeneous modes and uses them to choose an exact stationary reference.
Their constant leading terms arise from affine coordinate changes;
their normalization may also add decaying Kerr and gauge directions.
They therefore do not represent ten additional physical Kerr parameters.

\begin{proposition}[Exact stationary families and supplements]
\label{st:families}
For small $p=(m,v)\in\mathbb R^4\times\mathbb R^3$ there is
a smooth family $G_p$ of exact diffeomorphic stationary Kerr
metrics with $G_0=b$, tangent directions $H,\mathcal L_Wb$,
and symbol asymptote
\begin{equation}\label{st:affine-asymptote}
 G_p-\eta_v=O(r^{-1}),\qquad
 \eta_v=-(1-|v|^2)dt^2-2v_i\,dy_i\,dt+|dy|^2.
\end{equation}
Let $L_p(0)$ denote the stationary linear reduced operator at
metric and target $G_p$.  It has seven independent decaying
modes and ten further independent stationary homogeneous
metric modes with constant tensor asymptotes.  There are no
additional dimensions at any growth power strictly between
$0$ and $1$.

There exist a fixed compact smooth seven-vector functional
$\ell$, smooth ten-vector functionals $\ell_C^p$, and a smooth
list $H_C(p)$ of the ten additional modes, satisfying
\[
 \ell(H,J)=I,\qquad
 \ell H_C(p)=0,\qquad
 \ell_C^p H_C(p)=I,
\]
with $\ell_C^p$ annihilating the seven decaying modes.
Their supports can be fixed outside a physical trapping
neighborhood.  There is a second smooth exact stationary Kerr
family $\overline G_{p,c}$, $c\in\mathbb R^{10}$ small, with
\begin{equation}\label{st:second-family}
 \overline G_{p,0}=G_p,\qquad
 \partial_c\overline G_{p,0}=H_C(p).
\end{equation}
The gauge calibration extends smoothly to both families.
The frozen gauge of $\overline G_{p,c}$ relative to $G_p$
is $O(c^2)$ on fixed compact sets.
\end{proposition}

\begin{proof}
Integrate the stationary correction flows in
Lemma~\ref{st:profiles}.  For the velocity directions use
the affine leading chart
\begin{equation}\label{st:affine-chart}
 t_K=\sqrt{1-|v|^2}\,t+
              \frac{v\cdot y}{\sqrt{1-|v|^2}},
 \qquad
 y_K=\left(I+\frac{vv^t}{1-|v|^2}\right)^{1/2}y.
\end{equation}
Pullback of the flat metric is $\eta_v$.  Nearby raw Kerr
metrics are first expressed in smooth ingoing coordinates at
their outer horizons, bent on a larger domain, and corrected
to Boyer--Lindquist conventions towards infinity before
restriction.  Since the center has nonzero spin, its direction
can be rotated smoothly relative to the fixed Cartesian axes.
The correcting maps have stationary spatial components and
time component equal to a nonzero constant times $t$ plus a
spatial offset.  Their vector-field derivatives retain the
tight $r^{-1}$ bounds.  This constructs $G_p$ and
\eqref{st:affine-asymptote}.  Apply the smooth calibration at
$G_p$.

The supplemented stationary and gauge-potential estimates give
seven gauge-corrected decaying modes and no more.  Physical
variations are taken with the leading affine chart held fixed.
At vector growth just above constants the correction equations
are surjective, as in Lemma~\ref{st:profiles}.
Ten independent constant tensor asymptotes are obtained from
pure affine gauge at infinity, stationary as metric variations,
including a scaling of the stationary generator; correct them
by stationary $O(r^{0+})$ vectors.  The flat indicial gap,
unchanged by the constant anisotropic chart, shows that these
are exactly the ten extra dimensions below growth power $1$.
Uniform supplemented estimates at each fixed high order and
difference quotients of the differentiated stationary equations
prove smooth parameter dependence, with only finitely many
extra orders at each step.

Normalize their representatives by $\ell H_C(p)=0$.
For the probes, choose fixed tests separating the limiting
spans on a far exterior region.  The stationary estimates with
compact error, the dimension counts, and equation convergence
prevent any normalized nearby mode from being annihilated by
all of them.  Smooth parameter-dependent linear combinations
therefore produce $\ell_C^p$ with fixed supports.
These supports will lie in the unsmoothed region of the causal
construction.  Integrating the Kerr and diffeomorphism
directions that define $H_C(p)$ gives
$\overline G_{p,c}$.  Its slightly looser symbol decay after
its affine asymptote is harmless.  Continue the calibration
at this family.  Since its tangent satisfies the linear
stationary gauge, its frozen gauge residual is quadratic in
$c$.  All horizon, slicing, trapping, and conic-support
conditions retain strict margins after shrinking the parameter
range.
\end{proof}

The two families will serve different parts of the finite-time
argument.  Causal modulation will control the seven decaying modes,
while the ten measured coefficients $c$ select the instantaneous
reference $\overline G_{p,c}$ in the sliced estimate.  The difference
from this reference, rather than the entire perturbation, will have
the integrable low norms needed in the energy comparison.

For later energy estimates the frozen transport operator is now
unambiguously
\begin{equation}\label{st:frozen-operator}
 \mathbf L_{p,c}
 =\operatorname{Ric}'_{\overline G_{p,c}}
 +(\mathcal D_{\overline G_{p,c}}+C(p,c))
       \bigl(D\Upsilon_{\mathrm{wave}}|_{\overline G_{p,c};G_p}
                               +E_{\mathrm g}(p,c)\bigr).
\end{equation}
It does not include the derivative of $\mathcal D_g$ applied
to the residual gauge.  That omitted coefficient is
$O(c^2)$ compactly and will be treated in the nonlinear
comparison.  The construction order has been $b$, then $G_p$,
then $\overline G_{p,c}$.

\begin{lemma}[Adjoint tails]\label{st:adjoint-tails}
Let $\mathcal K_p^*$ be the supported bilinear-spatial-transpose
kernel dual to the decaying stationary problem.  It is
seven-dimensional, locally of order $\dot H^{-4}$, and
conormal $O(r^{-2+})$ at infinity.  It has locally uniform
weak bases with uniform symbol bounds sufficiently far out.
If $H_0$ is a smooth extendible stationary tensor of symbol
growth $O(r^{0+})$, then
\begin{equation}\label{st:adjoint-pairing}
 \langle v,L_p(0)H_0\rangle=0,\qquad v\in\mathcal K_p^*.
\end{equation}
\end{lemma}

\begin{proof}
Index and kernel dimension give the dimension. Apply Green's identity
to the constant-coefficient flat, possibly anisotropic, normal operator.  Indicial regularity
initially permits order-$(-1)$ fundamental solutions.  Pairing
with the ten independent constant-asymptote homogeneous modes of
\Cref{st:families} kills their leading coefficients.  The supported convention
supplies no inner boundary contribution, so the remaining
adjoint tail is $O(r^{-2+})$.  The uniform stationary estimates
and exterior indicial regularity give the asserted weak bases
and symbol bounds.

For $v=O(r^{-2+\epsilon'})$ and $H_0=O(r^\epsilon)$ the outer
Green boundary terms are $O(r^{-1+\epsilon+\epsilon'})$.
Choose $\epsilon+\epsilon'<1$ and let the sphere tend to
infinity.  The inner term vanishes by supported--extendible
duality.  This proves \eqref{st:adjoint-pairing}; the extra
adjoint tail power, rather than dual Fredholm information at
one weight alone, is what permits these pairings.
\end{proof}

\subsection{The causal target and reduced evolution}\label{st:reduction}

We now construct the time-dependent target and causal feedback that
will control the stationary zero modes. The radius-dependent averages
preserve the asymptotic symbol class; the filters control the parameter
jets; and the compensator restores the desired stationary linearization.
After proving their defect bounds, we solve the reduced evolution and
verify propagation of the vacuum constraints. Write $p(t)=(m(t),v(t))$,
let $a(t)=p'(t)$, and let $\lambda'(t)=v(t)$.
Extend $p$ and $\lambda$ by zero for $t\leq0$.
The co-moving coordinates are $(t,y)$ and the laboratory
coordinates are $(t,z)$, where
\[
 z=y-\lambda(t).
\]
The symbol $a$ in this section denotes the parameter derivative,
not the fixed Kerr rotation parameter.

Choose a smooth function $d(x)\geq0$, zero on a large fixed
ball containing both variables of every compact kernel and all
probes, and equal to $|x|$ near infinity.  Choose a real smooth
signed kernel $k_0$ compactly supported in $(0,1/10)$ with
\[
 \int k_0(\sigma)\,d\sigma=1,\qquad
 \int \sigma k_0(\sigma)\,d\sigma=0.
\]
For a zero-past function $f$ define its radius-dependent causal
average by
\begin{equation}\label{st:averaging}
 f_s(t,x)=\int k_0(\sigma)f(t-\sigma d(x))\,d\sigma .
\end{equation}
The implicit equation
\begin{equation}\label{st:implicit-position}
 x=z+\lambda_s(t,x)
\end{equation}
is a contraction for sufficiently small $\sup|v|$.
It gives $x=y$ in the unsmoothed region, and $1+|x|$ is
comparable to $1+|y|$.

Define the target by evaluating the family coefficients in the
formal coframe
\begin{equation}\label{st:causal-target}
 B=G_{p_s(t,x)}(x)
       \quad\hbox{in the coframe }(dt,dz+v_s(t,x)\,dt).
\end{equation}
The zero first moment has a useful exact consequence.  If $p$ is
constant and $\lambda$ is affine, with $\lambda'=v$, throughout the
interval from the earliest sampled time to the current time $t$, then
$p_s=p(t)$, $\lambda_s=\lambda(t)$, and
$x=y$.  On a region where these identities hold, $B$ is exactly $G_p$
in the co-moving chart, expressed in the laboratory chart by the
coframe $dy=dz+v\,dt$.  Thus the construction changes an exact
stationary metric only where the parameter history changes; the
following defect estimates quantify that change.

For the evolving metric $g$, set
\begin{equation}\label{st:unknowns}
 h=g-B,\qquad
 u=\partial_t h_{\mathrm{co}},\qquad
 D=\partial_{t,\mathrm{lab}}-v_s\cdot\partial_z,\qquad
 \psi=D h_{\mathrm{lab}}.
\end{equation}
Subscripts specify the coordinate components and derivatives.

Fix a smooth startup cutoff $\chi_{\mathrm{st}}$, zero near
zero and one after time one.  The feedback is
\begin{equation}\label{st:feedback}
 \begin{split}
 c&=\ell_C^p h_{\mathrm{co}},\\
 a&=F(\partial_t)(\chi_{\mathrm{st}}\ell h_{\mathrm{co}}),
 \qquad
 F=\operatorname{diag}(F_m I_4,F_v I_3),\\
 F_m(\zeta)&=\frac{\sigma_0}{(1+\zeta)^P},\qquad
 F_v(\zeta)=\frac{2\sigma_0\zeta+\sigma_0^2}{(1+\zeta)^P}.
 \end{split}
\end{equation}
Here $P\geq5$ is a fixed integer, increased below to dominate
a fixed finite list of frequency losses, and $\sigma_0>0$
is sufficiently small.  Both filters are retarded.

\begin{lemma}[Filter regularity]\label{st:filter}
For each $j\geq0$, the $j$th derivative of $a$ is controlled
by derivatives of its input through order
$\max(0,j-P+1)$ in polynomially weighted $L^2_t$ and
$L^\infty_t$, also on finite intervals.  For $j\geq1$ one
may first differentiate the input and use $\ell u$, with
the resulting startup term charged to initial data.
\end{lemma}

\begin{proof}
Put $d_n=\max(0,n-P+1)$ and let $q$ be the smooth zero-past input.
Let \(G_0=\delta_0\) and, for \(j\ge1\), let
\[
 G_j(t)={\bf1}_{t\ge0}\,e^{-t}\frac{t^{j-1}}{(j-1)!}.
\]
Write \(K_F\) for the diagonal convolution kernel of \(F\). Its mass/spin kernel is \(\sigma_0G_P\), and its velocity kernel is
\[
 2\sigma_0G_{P-1}+(\sigma_0^2-2\sigma_0)G_P,
\]
since \(\zeta(1+\zeta)^{-P}=(1+\zeta)^{-(P-1)}-(1+\zeta)^{-P}\). The distributional identity
\(\partial_tG_j=G_{j-1}-G_j\) shows that every derivative through order \(P-1\) of these kernels is a finite measure, including its possible endpoint delta. All its polynomial moments are finite. Put \(r_n=\min\{n,P-1\}\). Differentiating the convolution gives
\[
 a^{(n)}=(\partial_t^{r_n}K_F)*q^{(d_n)}.
\]
There are no startup delta terms from \(q\), since its zero extension is smooth. For \(0\le u\le t\),
\[
 \frac{(T_0+t)^{\alpha/2}}{(T_0+t-u)^{\alpha/2}}
 \le (1+u)^{\alpha/2}.
\]
Weighted Young's inequality for this finite measure proves the asserted
weighted bounds, on finite intervals by causality. For the input
$q=\chi_{\mathrm{st}}\ell h_{\mathrm{co}}$, differentiating once uses
$\ell u$ and a fixed startup term. The primitive relation
$p^{(j)}=a^{(j-1)}$ gives no strong weight for $p$ or $\lambda$.
\end{proof}

Use
\begin{equation}\label{st:full-gauge}
 \Upsilon=\Upsilon_{\mathrm{wave}}(g;B)
                   +E_{\mathrm g}(p,c)h+\Theta
\end{equation}
in \eqref{st:reduced-normalization}, with partner $C(p,c)$.
The spatial additions act in the co-moving chart.
To define the compensator, contract the profile indices and put
\begin{equation}\label{st:compensator}
 Y=H m_s+J\lambda_s+H^1v_s
               +2(V-e)^\flat\odot d_x\lambda_s .
\end{equation}
On the unsmoothed region, containing all kernels with a buffer,
set $\Theta=\mathcal B Y$, with the base coefficients of
$\mathcal B$.  Outside, compute the same local expression in
independent coordinates $(t,x)$, then evaluate at the implicit
$x$ and pull its components into the formal coframe of
\eqref{st:causal-target}.  The definitions coincide where $d=0$.

\begin{lemma}[Joint linearization and column bounds]\label{st:columns}
Linearizing the joint reduced equation at zero parameters and
$g=b$ gives
\begin{equation}\label{st:joint-linearization}
 Lh+Q(\partial_t)a.
\end{equation}
For
\[
 K_d(y,\zeta)=\int k_0(\sigma)e^{-\zeta\sigma d(y)}\,d\sigma
\]
the seven columns of $Q$ are
\begin{equation}\label{st:Q-definition}
 \begin{split}
 Y_m&=H K_d,\\
 Y_v&=(J+\zeta H^1)K_d
                   +2(V-e)^\flat\odot d_yK_d,\\
 Q(\zeta)&=
   \left(\frac{L(\zeta)Y_m}{\zeta},
         \frac{L(\zeta)Y_v}{\zeta^2}\right).
 \end{split}
\end{equation}
The apparent poles are removable.  On every fixed compact set,
$Q(\partial_t)$ has only bounded-length causal memories,
including distributional derivatives, and its fixed spatial
derivative costs are polynomial in frequency.  For every
fixed $N$,
\begin{equation}\label{st:Q-tail}
 Q(\zeta)=O_N\bigl(r_y^{-2}(1+|\zeta|r_y)^{-N}\bigr),
 \qquad r_y=|y|,\quad \operatorname{Re}\zeta\geq0,
\end{equation}
with the corresponding symbol derivatives near infinity.
\end{lemma}

\begin{proof}
At zero parameters the variation of $B$ in the fixed co-moving
chart, plus the Lie derivative along
$e\lambda+(V-e)\lambda_s$, is precisely $Y$.
The location variation supplies $\lambda_s-\lambda$ and the
formal-coframe variation supplies $v_s-v$.
Since $b$ is vacuum, adding this Lie derivative does not alter
the linearized Ricci tensor.  The definition of $\Theta$
therefore gives \eqref{st:joint-linearization}, with the columns
\eqref{st:Q-definition}.

The stationary identities and the chain identity
\eqref{st:profile-identities} cancel the zeroth and first
Taylor terms required by the displayed denominators.
The zeroth and first moments of $k_0$ give
$K_d=1+O(|\zeta|^2r_y^2)$ at small scaled frequency.
Thus both poles are removable; in the time domain their
antiderivatives have bounded memory on bounded spatial sets.
At large scaled frequency integrate by parts in $\sigma$.
Its compact support is strictly at positive delays, so the
bound is uniform for $\operatorname{Re}\zeta\geq0$.
Together with the tight profile orders
$H,H^1,V-e=O(r_y^{-1})$ and $J=O(r_y^{-2})$, these
calculations prove \eqref{st:Q-tail}.  Spatial derivatives
on a fixed annulus give only fixed polynomial frequency costs.
\end{proof}

\begin{lemma}[Nonlinear target defect]\label{st:target-defect}
Suppose $\int|a|$ and $\sup|v|$ are sufficiently small.
For comparable large radii $r\sim|x|$, set
\[
 \widetilde a(t,r)
       =Cr^{-1}\int_{t-Cr}^{t}|a(\tau)|\,d\tau,
\]
with the positive majorant enlarged as needed at comparable
radii.  In laboratory components,
\begin{equation}\label{st:target-defect-orders}
 \Theta=O(r^{-1}\widetilde a),\qquad
 f_B:=\operatorname{Ric}(B)
          +(\mathcal D_B+C(p,c))\Theta
                =O(r^{-2}\widetilde a),\qquad
 Df_B=O(r^{-3}\widetilde a).
\end{equation}
These hold with each fixed number of scaled laboratory
derivatives.  Also $DB=O(r^{-1}\widetilde a)$ with symbol
derivatives.  On fixed transition and inner sets the analogous
defects use bounded-length parameter jets and vanish at
exact freezing with affine position.
\end{lemma}

\begin{proof}
Work first where $d(x)=r=|x|$, with the enlarged positive majorant
$\widetilde a=Cr^{-1}\int_{t-Cr}^t|a(\tau)|\,d\tau$.
Assume the finite-slab bootstrap $\int|a|$ is small, so $r\widetilde a\ll1$.
Integration by parts in the smooth positive-delay kernel, with no endpoint terms, gives
\begin{equation}\label{st:smoothed-jets}
 \begin{aligned}
 |\partial_t^j\partial_x^\alpha p_s|
   &\le C_{j,\alpha}r^{1-j-|\alpha|}\widetilde a,
       &&j+|\alpha|\ge1,\\
 |\partial_t^j\partial_x^\alpha J_s|
   &\le C_{j,\alpha}r^{1-j-|\alpha|}\widetilde a,
       &&j,|\alpha|\ge0,\qquad J_s=\partial_x\lambda_s.
 \end{aligned}
\end{equation}
For the second line the zero first moment removes the constant-velocity term: \(J_s=O(r\widetilde a)\), while \(\partial_tJ_s=\partial_xv_s=O(\widetilde a)\). Each differentiated kernel is estimated separately by the positive majorant; the argument never differentiates \(\widetilde a\).

Put \(S_x=(I-J_s)^{-1}\). Implicit differentiation yields, in the independent variables \((t,x)\),
\begin{equation}\label{st:target-frame}
 \partial_{z_i}=(S_x)^j{}_i\partial_{x_j},\qquad
 D:=\partial_{t,\mathrm{lab}}-v_s\cdot\partial_z=\partial_t|_x,
 \qquad [D,\partial_{z_i}]
       =(\partial_{z_i}v_s^j)\partial_{z_j}.
\end{equation}
Thus \(S_x-I=O(r\widetilde a)\), \(DS_x=O(\widetilde a)\), and \(\partial_xS_x=O(\widetilde a)\), with one inverse radius for each further ordinary derivative. Products of acceleration averages obey the same bounds since \(r\widetilde a\ll1\).

\paragraph{The metric defect.}
The formal coframe substitution cancels the affine flat asymptote exactly: substituting \(dy=dz+v\,dt\) in \(\eta_v\) gives \(\eta\). Hence in lab components
\[
 B=\eta+\mathfrak h(p_s,x),\qquad
 \partial_x^\alpha\partial_p^\beta\mathfrak h=O(r^{-1-|\alpha|}).
\]
Here \(\mathfrak h\) is the smooth family of lab-component metric remainders. At a point freeze \(p_s,v_s\), and use the constant-parameter affine-worldline Kerr metric with position \(x\). Its first and second lab derivatives differ from those of \(B\) by
\(O(r^{-1}\widetilde a)\) and \(O(r^{-2}\widetilde a)\), respectively. Indeed the second derivative expansion consists of
\[
 \mathfrak h_{xx}CC,\quad \mathfrak h_x\partial C,\quad
 \mathfrak h_{xp}C\partial p_s,\quad \mathfrak h_{pp}(\partial p_s)^2,\quad
 \mathfrak h_p\partial^2p_s,
\]
with both mixed terms understood, where \(C\) is \(S_x\) or \(v_sS_x\). The coordinate-coefficient difference is \(O(r\widetilde a)\), \(\partial C=O(\widetilde a)\), and \(\partial^2p_s=O(r^{-1}\widetilde a)\). These give the claimed orders. The frozen metric is vacuum, so smooth expansion of Ricci gives \(\operatorname{Ric}(B)=O(r^{-2}\widetilde a)\).

For the additional \(D\) derivative, \(DB=O(r^{-1}\widetilde a)\) with two further lab symbol derivatives by \eqref{st:smoothed-jets}. For any two lab coordinate derivatives, the exact commutator identity is
\[
\begin{split}
 D\partial_\alpha\partial_\beta B
 ={}&\partial_\alpha\partial_\beta(DB)
 +(\partial_\alpha v_s^j)\partial_j\partial_\beta B
 +(\partial_\beta v_s^j)\partial_\alpha\partial_j B\\
 &+(\partial_\alpha\partial_\beta v_s^j)\partial_jB.
\end{split}
\]
Here \(\partial v_s=O(\widetilde a)\), \(\partial^2v_s=O(r^{-1}\widetilde a)\), \(\partial B=O(r^{-2})\), and \(\partial^2B=O(r^{-3})\). Every displayed term is therefore \(O(r^{-3}\widetilde a)\). Differentiating the coefficients in Ricci adds only terms of that order or smaller. Thus \(D\operatorname{Ric}(B)=O(r^{-3}\widetilde a)\).

\paragraph{The compensator and its derivative.}
At infinity the gauge is local; write it as
\(\mathcal B=\beta^\mu(x)\partial_\mu+\beta_0(x)\), and put \(U_i=V_i-e_i\). The identities
\[
 \mathcal BH=\mathcal BJ=0,\qquad
 \mathcal BH^1=-[\mathcal B,t]J,\qquad
 \partial_t\lambda_s=v_s
\]
reduce \(\mathcal BY\) to sums of
\[
 [\mathcal B,m_s]H,\qquad
 \beta^iJ\,\partial_{x_i}\lambda_s+[\mathcal B,v_s]H^1,
 \qquad
 \mathcal B(2U^\flat\odot d_x\lambda_s).
\]
In particular no undifferentiated position or velocity remains. The profile orders
\(H,H^1,U=O(r^{-1})\), \(J=O(r^{-2})\) and
\eqref{st:smoothed-jets} bound each expression by
\(O(r^{-1}\widetilde a)\). A \(D=\partial_t|_x\) hit on a memory gains one inverse radius; for example it sends
\(\partial_x\lambda_s=O(r\widetilde a)\) to
\(\partial_xv_s=O(\widetilde a)\). The implicit-coordinate and coframe changes contribute only \(S_x,v_s\) and their controlled derivatives. A derivative on a coframe coefficient produces
\(r^{-1}\widetilde a^2\lesssim r^{-2}\widetilde a\) at the \(\Theta\) level and
\(r^{-2}\widetilde a^2\lesssim r^{-3}\widetilde a\) after one further derivative. Together with \eqref{st:target-frame}, this proves
\[
 \Theta=O(r^{-1}\widetilde a),\quad
 \partial\Theta,D\Theta=O(r^{-2}\widetilde a),\quad
 D\partial\Theta=O(r^{-3}\widetilde a).
\]
Christoffel terms in \(\mathcal D_B\Theta\) have smaller orders, and \(C\) vanishes here. We obtain the required defect bounds
\begin{equation}\label{st:full-defect-identity}
 \Theta=O(r^{-1}\widetilde a),\qquad
 f_B:=\operatorname{Ric}(B)+(\mathcal D_B+C)\Theta
       =O(r^{-2}\widetilde a),\qquad
 Df_B=O(r^{-3}\widetilde a).
\end{equation}
The calculation allows every fixed number of scaled lab derivatives, with constants depending on that number. Each use may enlarge the comparable averaging interval. At the far-current application, lab-time derivatives are taken before the outgoing commutations; their coefficients contain only smoothed parameters. No unsmoothed acceleration is created by the calculation above.

On the fixed transition annulus \(d\) may vanish, so an inverse-\(d\) estimate is neither needed nor asserted. There the defect is a smooth function of finitely many parameter jets and bounded-length memories. In the unsmoothed region it also contains the fixed compact pseudodifferential operations. The exact identities above make it vanish whenever the relevant accelerations vanish. The filter lemma supplies every fixed required jet list, and the compact estimates treat these transition terms with fixed-radius constants.

\end{proof}

The target, feedback, and compensator now have the spatial decay and
finite-jet control needed by the evolution. We next show that their
retarded dependence closes the local Cauchy problem and that vacuum
initialization gives the vacuum equation throughout the resulting slab.

For this local construction, the base wave-pair norm is
\[
 \mathcal E_{10}(t)=\|h_{\rm co}(t)\|_{H^{10}_{\rm uloc}}
                 +\|\partial_t h_{\rm co}(t)\|_{H^9_{\rm uloc}}.
\]
Here the uniformly local norms use a fixed buffered region containing
all compact kernels and regular end charts of fixed size. The contraction
norm is the wave-pair norm one order lower, $H^9_{\rm uloc}\times
H^8_{\rm uloc}$. These are sufficient fixed orders, not a minimal
regularity assertion.

\begin{proposition}[Local evolution and vacuum initialization]
\label{st:local-evolution}
The reduced equation
\eqref{st:reduced-normalization},
\eqref{st:full-gauge}, with retarded feedback
\eqref{st:feedback}, has smooth local evolution for smooth
initial wave data small in $\mathcal E_{10}$.
It continues on finite time intervals while $\mathcal E_{10}$ remains
bounded, the parameters stay in their chosen small range, the normalized
acceleration and implicit-coordinate inverses stay uniformly controlled,
and the spacelike and strict outflow margins remain positive.
Smooth higher orders propagate on bounded time intervals;
these higher seminorms need only be finite.
For smooth vacuum constraint data, the initial lapse and shift
derivatives can be chosen so that $\Upsilon$ and its first
normal derivative vanish initially.  The resulting evolution
is vacuum.
\end{proposition}

\begin{proof}
After pointwise normalization of the wave principal coefficient,
the extra normal-acceleration coefficient has compact interior
spatial order at most $-1$.  Choose it small on a low norm and
on the finite base Sobolev list needed for iteration.
Its identity-plus-correction inverse then extends to every
smooth order by the derivative gain and Sobolev multiplication.
The identity-minus-inverse error gains one spatial derivative;
apply it to $\partial_y^2h$ and $\partial_y\partial_th$.
The solved equation therefore retains its differential strictly
hyperbolic principal part, and every nonlocal correction is
at most first order on the wave jets, with tame bounds.

The compact kernels are bounded in a high spatial wave-pair
energy and are Lipschitz one energy order lower on bounded
sets.  Metric multiplication inside their negative-order gains
uses ordinary spatial Sobolev multiplication.  Before spatial
commutation, the equation contains at most $\partial_tc$
from the parameter-dependent additions, a projection of the
two available wave jets rather than a new unsolved acceleration.
At any fixed base iteration order choose $P$ with surplus so
that spatial differentiation of bounded-radius histories uses
only available projection jets.  The implicit-position
inverses use $I-\partial_x\lambda_s$ uniformly close to the
identity, and the filters are stable zero-past Volterra
variables, equivalently a finite system of retarded ordinary
differential equations.

On a short slab integrate the high-energy estimate and the
one-order-lower contraction estimate, using suprema of
earlier wave energies for the history terms.  Work on regions
containing the compact kernels with buffers and use uniformly
local wave estimates outside.  The inner boundary has strictly
exiting flux. Conic localization preserves the spatial mapping properties and
the time-local or retarded dependence. Both variables of every compact
spatial kernel stay within the buffered interaction region, away from
the outflow boundaries. This does not assert metric finite propagation
inside that region; later geometric attachment uses ordinary local
wave-map equations.
At higher orders, spatial commutation through source order
$s-1$ uses at most $a^{(s+3)}$ in targets and compensators.
The filter gain charges these to at most $s-1$ total-order
projection jets when $P\geq5$.
Successive normal solves and wave-energy induction, or smooth
regularization with these tame estimates, propagate smoothness
without demanding smallness of the negative-order coefficient
on every high Sobolev space.  Outside the compact interaction
region, evolution for fixed histories is local.  Annular
rescaling, together with the fact that sufficiently distant
memories sample times before startup, propagates all symbol
orders on bounded time slabs.

For constraint data, imposing $\Upsilon=0$ determines the
initial lapse and shift derivatives by the same small normal
coefficient solve.  The normal Einstein constraints in the
reduced equation then imply $\partial_t\Upsilon=0$ initially.
The contracted Bianchi identity gives a homogeneous wave
system for $\Upsilon$; the compact terms from $C$ have the
same lower-order type.  Uniqueness yields $\Upsilon=0$ and
$\operatorname{Ric}(g)=0$.  The startup cutoff on feedback
ensures that the initial wave jets through the finite base
budget are determined by the corresponding geometric data
orders.  None of these arguments requires stationary time
to be timelike on the kernel supports.
\end{proof}

\subsection{Augmented inversion and retarded tests}

The unaugmented operator has stationary modes.  The feedback removes
their obstruction without assuming a uniform unaugmented inverse
at zero.  For a large radius $R$, choose a smooth scaled cutoff
$\chi_R(y)$, equal to one for $r_y\leq R$ and zero for
$r_y\geq2R$, and define
\begin{equation}\label{st:augmented-operator}
 M_R(\zeta)=L(\zeta)+\chi_R Q(\zeta)F(\zeta)\ell.
\end{equation}
We first identify the zero-frequency coupling.

\begin{lemma}[Zero-frequency coupling]\label{st:zero-coupling}
For sufficiently large $\Lambda$, the map
\[
 Q(0)^t:\mathcal K_0^*\longrightarrow\mathbb C^7
\]
has full rank.
\end{lemma}

\begin{proof}
Expand the columns in outgoing time $t_*=t-\phi$, equivalently
multiplying the profiles by $e^{\zeta\phi}$.  In pairing their
Taylor coefficients with $\mathcal K_0^*$, discard terms
$L(0)H_0$ with $H_0=O(r_y^{0+})$ by
Lemma~\ref{st:adjoint-tails}.  This operation is essential for
the coefficient class of the eventual polynomial mode.
For a Kerr column the discarded coefficient includes $\phi H$.
For a translation column, the constant and linear coefficients
of the conjugated profile are $J$ and
$Z_1=H^1+\phi J=O(r^{-1+})$; its quadratic coefficient includes
$\phi H^1+\phi^2J/2$, together with the averaging and coordinate
profile terms.  The latter coefficients are $O(r^{0+})$ and
contribute only exact stationary terms $L(0)H_0$ at the order
under consideration.  They are discarded in the adjoint pairing,
not retained as constant coefficients of a polynomial to which
Proposition~\ref{st:hhv} is applied.  The residuals at the relevant
orders are exactly the equation errors of
\begin{equation}\label{st:polynomial-tests}
 t_*H,\qquad
 \tfrac12t_*^2J+t_*(H^1+\phi J).
\end{equation}
These errors are stationary and $O(r_y^{-3+})$.  To see the
extra decay, the outgoing linear Taylor coefficient on the end
is a constant multiple of $\partial_r+1/r$, modulo
$O(r^{-1})\partial_y+O(r^{-2})$, and annihilates the leading
$r^{-1}$ power.  The outgoing quadratic coefficient is
$O(r^{-1})$ or better.  These are the flat and Kerr asymptotic
computations using precisely the
orders of $b$, $\phi$, and the profiles established above.

At the local-gauge limit, if the resulting seven-dimensional
coupling were singular, a nonzero combination of these residuals
would pair to zero with the entire adjoint kernel.  The improved
$O(r_y^{-3+})$ residual could then be solved away in a decaying
$O(r_y^{-1+})$ space.  Add only this decaying stationary solution
to the corresponding combination in \eqref{st:polynomial-tests}.
Every coefficient of the resulting outgoing-time polynomial
then lies in the decaying class of Proposition~\ref{st:hhv};
none of the discarded $O(r^{0+})$ coefficients occurs in it.
If its quadratic coefficient were
nonzero, \eqref{st:polynomial-tests} would give a forbidden
degree-two outgoing polynomial mode.  Otherwise a nonzero
linear combination of the four Kerr directions would give a
degree-one mode with nontranslation leading coefficient.
Both contradict Proposition~\ref{st:hhv}.
It is essential here to discard exact stationary terms before
invoking solvability for the improved residual.

The added gauge preserves all identities used in these Taylor
expansions.  Uniform stationary compactness and the
$O(r_y^{-2+})$ adjoint tails give continuity of this finite
coupling towards the local-gauge limit.  Thus its rank persists
for sufficiently large $\Lambda$.
\end{proof}

\begin{proposition}[Retarded stationary tests]\label{st:retarded-tests}
Fix a compact test support and a finite number of required
angular derivatives.  There are a sufficiently high fixed
test Sobolev order $q_*$ and a finite polynomial degree $N_*$
with the following property.  Choose the filter order $P$
sufficiently large and $\sigma_0>0$ sufficiently small, then
take $\Lambda,R$ sufficiently large in the order specified
below.  For $0<\operatorname{Re}\zeta\leq D_*$ on that strip
and a compactly supported test $f$ of order $q_*$, the spatial
transpose equation
\[
 M_R(\zeta)^t v_*=f
\]
has a solution satisfying
\begin{align}
 \|v_*\|_{\dot H^{-4}(K)}
 &\leq C R^{3\iota}\langle\zeta\rangle^{N_*}
                       \|f\|_{H^{q_*}},
 \label{st:test-compact}\\
 \|(v_*,r_y(\partial_{r_y}+\zeta)v_*)\|_{L^2_\omega}
 &\leq C r_y^{-1+\iota}R^{3\iota}
                \langle\zeta\rangle^{N_*}\|f\|_{H^{q_*}}
       \qquad(r_y\gg1).
 \label{st:test-radiation}
\end{align}
Here $K$ is any prescribed fixed compact reach, and the second
bound includes the prescribed angular derivatives.
The bounds are uniform down to
$\operatorname{Re}\zeta\downarrow0$ on a fixed strip reaching
a forward Laplace line.  The slack $\iota>0$ can be arbitrarily
small without making $N_*$ diverge.  The associated tested
operators are retarded.
\end{proposition}

\begin{proof}
We give the zero, intermediate, and high frequency arguments
separately.

\emph{Near zero.}
Put
\[
 \delta_f=F(\zeta)^tQ(\zeta)^t\chi_Rv_*,
 \qquad
 L(\zeta)^t v_*=f-\ell^t\delta_f.
\]
Use the spaces $\mathcal A,\mathcal S$ of
Lemma~\ref{st:transpose}, and supplement their compact-error
estimate by a compact measurement $\Pi$ that separates
$\mathcal K_0^*$.  It controls both
$e^{\zeta\phi}v_*$ and $\delta_f$ with the extra forcing
$\ell^t\delta_f$ placed on the left.  Indeed in a normalized
contradicting sequence, the limiting stationary equation paired
with $H,J$ gives $\delta_f=0$, because $\ell(H,J)=I$.
The measurement $\Pi$ then kills its remaining stationary
kernel element.

Choose $v_0\in\mathcal K_0^*$ matching
$\Pi(e^{\zeta\phi}v_*)$.  Its conjugated equation residual is
$O(|\zeta|)|v_0|$ in $\mathcal S$, by the adjoint
$r_y^{-2+}$ tail with smaller slack and the stationary end
bounds; here $|v_0|$ denotes its finite-dimensional coefficient
size.  The supplemented estimate gives
\begin{equation}\label{st:near-zero-comparison}
 \|e^{\zeta\phi}v_*-v_0\|_{\mathcal A}+|\delta_f|
 \leq C\bigl(\|f\|_{H^{q_*}}+|\zeta||v_0|\bigr).
\end{equation}
By the column bounds, pairing the physical difference
$e^{-\zeta\phi}(e^{\zeta\phi}v_*-v_0)$ against $\chi_RQ$
costs at most
\[
 C_\iota\min(R,|\zeta|^{-1})^{2\iota}
       \bigl(\|f\|_{H^{q_*}}+|\zeta||v_0|\bigr).
\]
On normalized $v_0$, the pairing of
$e^{-\zeta\phi}v_0$ with $\chi_RQ(\zeta)$ converges uniformly
to $Q(0)^t v_0$ as $R\to\infty$ and $\zeta\to0$.
Lemma~\ref{st:zero-coupling} and nonsingularity of $F(0)$
therefore absorb the
$O(|\zeta|^{1-2\iota})|v_0|$ term after choosing $R$ large
and a sufficiently small fixed zero-frequency neighborhood.
This gives the claimed compact bound and trivial augmented
transpose kernel there.  Lemma~\ref{st:transpose} supplies
index zero and the return to physical inversion.  Its stronger
near-zero source estimate is applied only to these graph
states, whose forcing is compactly supported.
At infinity b-Sobolev regularity for $e^{\zeta\phi}v_*$
gives \eqref{st:test-radiation} in this regime, using
$r_y(1-\phi')=O(1)$.

\emph{Closed nonzero bands.}
Use the estimates for $L^t$ first with the full column $Q$,
whose tail rapidly decreases, and then treat the removed
tail $(1-\chi_R)Q$ as small.  Spatial integration by parts
reduces the full finite-dimensional condition to
\begin{equation}\label{st:filter-matrix}
 I+[I_H,I_J+\zeta\,\ell H^1]\,
       \operatorname{diag}(F_m/\zeta,F_v/\zeta^2).
\end{equation}
Here $[I_H,I_J]$ is the identity divided into its four and
three mode columns.  Formula \eqref{st:filter-matrix}
uses $d=0$ on the probe supports and the exact stationary
and chain identities.

This matrix is invertible in the closed right half-plane for
sufficiently small $\sigma_0$.  It is a perturbation of the
identity when $|\zeta|/\sigma_0$ is large.  When
$w=\zeta/\sigma_0$ stays bounded, multiply its four Kerr
columns by $w$ and its three translation columns by $w^2$.
The limiting diagonal blocks are $1+w$ and $(1+w)^2$,
which have no zeros for $\operatorname{Re}w\geq0$.
The zero-frequency neighborhood has already been handled;
these two regimes cover the remaining closed bounded bands.
Taking $R$ large absorbs the omitted tail.

\emph{High frequency and radiation.}
The polynomial transpose losses from
Proposition~\ref{st:stationary-estimates}, the polynomial
column bounds on compact sets, and their rapid tail decrease
are absorbed by taking $P$ large.  Away from zero,
$L^t v_*$ is supported in a fixed compact set.
At the end we obtain the radiation estimate as follows.
Commute rotations in Cartesian components, with cutoff in an
elliptic annulus.  Their order-two coefficient errors are
$O(r_y^{-2})$ by rotation invariance of the mass correction.
The scattering estimate and exterior spatial ellipticity
control these unscaled derivatives and fixed-annulus errors
with polynomial cost.  For $w_*=r_yv_*$ this gives, with
the requisite rotations,
\[
 (w_*,\partial_{r_y}w_*)\in
     L^2(r_y^{-1-2\iota}dr_y\,d\omega)
\]
with the asserted right side.
In tortoise coordinate the source-free radial equation is
\begin{equation}\label{st:transpose-radial}
 (\partial_{r_*}-\zeta)(\partial_{r_*}+\zeta)w_*
 =O(r_y^{-2})\langle\zeta\rangle^2
       O\bigl(\partial_\omega^{\leq2}w_*,
                \partial_\omega^{\leq1}\partial_{r_*}w_*\bigr).
\end{equation}
This follows by dividing by the radial second-derivative
coefficient in Boyer--Lindquist Kerr.  For positive real part,
integrate the first factor from infinity; Cauchy--Schwarz
with the scattering weight costs $O(r_y^{-1+\iota})$.
To see the exact radial gain, if $w_*=r_yv_*$, the scattering
weight gives
\[
 \int_r^\infty \rho^{-2}|w_*|\,d\rho
 \le
 \left(\int_r^\infty\rho^{-1-2\iota}|w_*|^2\,d\rho\right)^{1/2}
 \left(\int_r^\infty\rho^{-3+2\iota}\,d\rho\right)^{1/2}
 \lesssim r^{-1+\iota}.
\]
The same estimate holds for the finitely many angular and radial jets
on the right side of \eqref{st:transpose-radial}. The first integration
has a kernel bounded by one; the second costs only a constant proportional
to $1/\iota$. Reducing $\iota$ changes constants but not the polynomial
frequency degree, because all differential orders were fixed beforehand.
Integrate the second factor outwards.  Returning to the
ordinary good derivative uses $r_*'=1+O(r_y^{-1})$.
This proves \eqref{st:test-radiation}, with all compact
modifications interior to the computation.

\emph{Choice of constants and causality.}
First fix the finite testing and differentiation orders and
their polynomial losses.  Choose $P$ and then $\sigma_0$.
A forward energy line can be fixed uniformly before finalizing
$\Lambda,R$: the acceleration coefficients are small, the
compact symbols are uniformly bounded, and $Q$ acting on
the filtered parameters is a bounded energy forcing by its
memory estimates.  The high and near-zero thresholds are
uniform for sufficiently large $\Lambda$.  Enlarge $\Lambda$
for the intervening band, and then take $R$ sufficiently
large.  This is an augmented argument using column identities,
rank, and transpose graphs; it does not require persistence
of unaugmented inversion uniformly to zero.

On the open strip, the physical graph domains are fixed
locally in $\zeta$: the end is parameter elliptic and compact
first-order changes are bounded on the graph.  The inverse is
analytic there.  Starting on the forward line, wave and
Volterra uniqueness identify the tested inverse as causal.
Contour displacement within the strip, using
\eqref{st:test-compact}--\eqref{st:test-radiation}, gives
the retarded tested operators down to the boundary of the
physical half-plane.
\end{proof}

\begin{remark}[Fixed base choices and higher regularity]
\label{st:order-quantifiers}
The choice of $\Lambda$ requires smallness on a finite base
Sobolev list, not on all orders simultaneously.  Once the
negative-order acceleration correction is invertible at a low
order, its derivative gain yields estimates at every fixed
higher order, with order-dependent constants.  The saturated
transport calibration is exact and independent of the energy
order.  For nearby $p$, the seven-dimensional kernel is
reconstructed using the invertible gauge-potential problem:
four physical variations and three translations give the
lower bound, while stationary compactness gives the matching
upper bound.  Preservation of the base translation chains is
needed for the base columns and coupling, not asserted as
an unchanged chain at every nearby parameter.

Finally the worst filter order is $1-P$, so
$a^{(s+3)}$ uses at most
$\max(0,s+4-P)\leq s-1$ projection derivatives when $P\geq5$.
This count is uniform in $s$.  Thus the fixed base construction
does not introduce a requirement of smallness at arbitrarily
high data orders.  The subsequent energy argument can use the
order-dependent stationary constants when proving smooth
higher bounds for each fixed datum.
\end{remark}
